% GENERATED FILE. Edit canonical lessons, then run npm run generate:books.
% canonical-source-hash: fnv1a64:a21516a7df43866c
% canonical-lessons: TA-C351-kurippitu, TA-C351-cuttikkattu, TA-C351-etutturai, TA-C351-atavatu, TA-C351-etuttukkattu

\chapter{Making a Point}
\label{ch:ta-making-a-point}
\hlchaptermodality{351}

\begin{chapteropening}
\textbf{By the end of this chapter:} \emph{I can mention, point out and explain what I mean.}

Complete the last lesson of chapter 351: I can mention, point out and explain what I mean.
\end{chapteropening}

\section[\texorpdfstring{\ta{குறிப்பிடு} (kuṟippiṭu)}{குறிப்பிடு (kuṟippiṭu)}]{\ta{குறிப்பிடு} (kuṟippiṭu) — to mention; to specify}

\label{lesson:TA-C351-kurippitu}



\begin{quote}
Before the new one: say the Tamil for an angle, then the Tamil for to express.
\end{quote}

\subsection*{You'll want to know: \ta{குறிப்பிடு}}
\textbf{\ta{குறிப்பிடு}} — \emph{kuṟippiṭu} — \textquotedblleft{}to mention; to specify\textquotedblright{}.

\begin{quote}
\ta{அவர்} \ta{தன்} \ta{கடிதத்தில்} \ta{இந்தப்} \ta{பிரச்சினையைக்} \ta{குறிப்பிட்டார்}.
\end{quote}

\begin{quote}
\emph{avar taṉ kaṭitattil intap piracciṉaiyaik kuṟippiṭṭār.}
\end{quote}

\begin{quote}
He mentioned this problem in his letter.
\end{quote}

Past (I …): \ta{குறிப்பிட்டேன்}, \emph{kuṟippiṭṭēṉ}.

What is mentioned takes -\ta{ஐ}. \ta{குறிப்பிட்ட}, specific: \ta{குறிப்பிட்ட} \ta{நேரம்}, a set time.

\begin{cousinweb}
\textbf{\ta{குறிப்பிடு}} (\emph{kuṟippiṭu}) is native: \emph{kuṟippu}, \textquotedblleft{}a mark, a sign\textquotedblright{}, belongs to \emph{kuṟi}, \textquotedblleft{}to aim at, to mark\textquotedblright{} (DEDR 1847), and \emph{iṭu}, \textquotedblleft{}to put\textquotedblright{} (DEDR 442), makes the verb: \textquotedblleft{}to put a mark on\textquotedblright{}.
\end{cousinweb}

\begin{grammarlens}[title={What you've built}]
One more word for giving reasons and opinions.
\end{grammarlens}

\subsection*{Your turn}
\begin{itemize}
\raggedright
  \item \emph{Say it:} \emph{kuṟippiṭu}
  \item \emph{Say it:} \emph{kuṟippiṭu}, once more
  \item \emph{Say it:} \emph{avar taṉ kaṭitattil intap piracciṉaiyaik kuṟippiṭṭār.} — He mentioned this problem in his letter.
  \item \emph{Recall:} say the Tamil for an angle, then the Tamil for to express, then say \emph{kuṟippiṭu} again
\end{itemize}

\begin{tcolorbox}[breakable,colback=teal!4,colframe=teal!35!black,title={Before you move on}]
What does \ta{குறிப்பிடு} mean? (\textquotedblleft{}To mention; to specify\textquotedblright{}.) Say it once more.
\end{tcolorbox}
\section[\texorpdfstring{\ta{சுட்டிக்காட்டு} (cuṭṭikkāṭṭu)}{சுட்டிக்காட்டு (cuṭṭikkāṭṭu)}]{\ta{சுட்டிக்காட்டு} (cuṭṭikkāṭṭu) — to indicate; to point out}

\label{lesson:TA-C351-cuttikkattu}



\begin{quote}
Before the new one: say the Tamil for to express, then the Tamil for to mention.
\end{quote}

\subsection*{You'll want to know: \ta{சுட்டிக்காட்டு}}
\textbf{\ta{சுட்டிக்காட்டு}} — \emph{cuṭṭikkāṭṭu} — \textquotedblleft{}to indicate; to point out\textquotedblright{}.

\begin{quote}
\ta{ஆசிரியர்} \ta{என்} \ta{தவறுகளைச்} \ta{சுட்டிக்காட்டினார்}.
\end{quote}

\begin{quote}
\emph{āciriyar eṉ tavaṟukaḷaic cuṭṭikkāṭṭiṉār.}
\end{quote}

\begin{quote}
The teacher pointed out my mistakes.
\end{quote}

Past (I …): \ta{சுட்டிக்காட்டினேன்}, \emph{cuṭṭikkāṭṭiṉēṉ}.

What is pointed out takes -\ta{ஐ}: \ta{தவறைச்} \ta{சுட்டிக்காட்டு}.

\begin{cousinweb}
\textbf{\ta{சுட்டிக்காட்டு}} (\emph{cuṭṭikkāṭṭu}) is native: \emph{cuṭṭi}, \textquotedblleft{}having indicated\textquotedblright{}, from \emph{cuṭṭu}, \textquotedblleft{}to point out, designate\textquotedblright{} (DEDR 2658), joined to \emph{kāṭṭu}, \textquotedblleft{}to show\textquotedblright{} (DEDR 1443).
\end{cousinweb}

\begin{grammarlens}[title={What you've built}]
One more word for giving reasons and opinions.
\end{grammarlens}

\subsection*{Your turn}
\begin{itemize}
\raggedright
  \item \emph{Say it:} \emph{cuṭṭikkāṭṭu}
  \item \emph{Say it:} \emph{cuṭṭikkāṭṭu}, once more
  \item \emph{Say it:} \emph{āciriyar eṉ tavaṟukaḷaic cuṭṭikkāṭṭiṉār.} — The teacher pointed out my mistakes.
  \item \emph{Recall:} say the Tamil for to express, then the Tamil for to mention, then say \emph{cuṭṭikkāṭṭu} again
\end{itemize}

\begin{tcolorbox}[breakable,colback=teal!4,colframe=teal!35!black,title={Before you move on}]
What does \ta{சுட்டிக்காட்டு} mean? (\textquotedblleft{}To indicate; to point out\textquotedblright{}.) Say it once more.
\end{tcolorbox}
\section[\texorpdfstring{\ta{எடுத்துரை} (eṭutturai)}{எடுத்துரை (eṭutturai)}]{\ta{எடுத்துரை} (eṭutturai) — to expound; to set out in full}

\label{lesson:TA-C351-etutturai}



\begin{quote}
Before the new one: say the Tamil for to mention, then the Tamil for to indicate.
\end{quote}

\subsection*{You'll want to know: \ta{எடுத்துரை}}
\textbf{\ta{எடுத்துரை}} — \emph{eṭutturai} — \textquotedblleft{}to expound; to set out in full\textquotedblright{}.

\begin{quote}
\ta{அவர்} \ta{தன்} \ta{காரணங்களை} \ta{எடுத்துரைத்தார்}.
\end{quote}

\begin{quote}
\emph{avar taṉ kāraṇaṅkaḷai eṭutturaittār.}
\end{quote}

\begin{quote}
He set out his reasons.
\end{quote}

Past (I …): \ta{எடுத்துரைத்தேன்}, \emph{eṭutturaittēṉ}.

A formal written word; the everyday phrase is \ta{எடுத்துச்} \ta{சொல்}.

\begin{cousinweb}
\textbf{\ta{எடுத்துரை}} (\emph{eṭutturai}) is native: \emph{eṭuttu}, \textquotedblleft{}having taken up\textquotedblright{}, from \emph{eṭu} (DEDR 851), joined to \emph{urai}, \textquotedblleft{}to speak, tell\textquotedblright{} (DEDR 648): to take up a point and speak it.
\end{cousinweb}

\begin{grammarlens}[title={What you've built}]
One more word for giving reasons and opinions.
\end{grammarlens}

\subsection*{Your turn}
\begin{itemize}
\raggedright
  \item \emph{Say it:} \emph{eṭutturai}
  \item \emph{Say it:} \emph{eṭutturai}, once more
  \item \emph{Say it:} \emph{avar taṉ kāraṇaṅkaḷai eṭutturaittār.} — He set out his reasons.
  \item \emph{Recall:} say the Tamil for to mention, then the Tamil for to indicate, then say \emph{eṭutturai} again
\end{itemize}

\begin{tcolorbox}[breakable,colback=teal!4,colframe=teal!35!black,title={Before you move on}]
What does \ta{எடுத்துரை} mean? (\textquotedblleft{}To expound; to set out in full\textquotedblright{}.) Say it once more.
\end{tcolorbox}
\section[\texorpdfstring{\ta{அதாவது} (atāvatu)}{அதாவது (atāvatu)}]{\ta{அதாவது} (atāvatu) — that is to say; in other words}

\label{lesson:TA-C351-atavatu}



\begin{quote}
Before the new one: say the Tamil for to indicate, then the Tamil for to expound.
\end{quote}

\subsection*{You'll want to know: \ta{அதாவது}}
\textbf{\ta{அதாவது}} — \emph{atāvatu} — \textquotedblleft{}that is to say; in other words\textquotedblright{}.

\begin{quote}
\ta{அவர்} \ta{நாளை} \ta{வருவார்}, \ta{அதாவது} \ta{திங்கள்கிழமை}.
\end{quote}

\begin{quote}
\emph{avar nāḷai varuvār, atāvatu tiṅkaḷkiḻamai.}
\end{quote}

\begin{quote}
He will come tomorrow, that is, on Monday.
\end{quote}

It introduces a restatement: \ta{அதாவது}, that is.

\begin{cousinweb}
\textbf{\ta{அதாவது}} (\emph{atāvatu}) is native: \emph{atu}, \textquotedblleft{}that\textquotedblright{} (DEDR 1), and \emph{āvatu}, \textquotedblleft{}that which is\textquotedblright{}, a verbal noun of \emph{āku}, \textquotedblleft{}to become\textquotedblright{} (DEDR 333): \textquotedblleft{}that which it is\textquotedblright{}.
\end{cousinweb}

\begin{grammarlens}[title={What you've built}]
One more word for giving reasons and opinions.
\end{grammarlens}

\subsection*{Your turn}
\begin{itemize}
\raggedright
  \item \emph{Say it:} \emph{atāvatu}
  \item \emph{Say it:} \emph{atāvatu}, once more
  \item \emph{Say it:} \emph{avar nāḷai varuvār, atāvatu tiṅkaḷkiḻamai.} — He will come tomorrow, that is, on Monday.
  \item \emph{Recall:} say the Tamil for to indicate, then the Tamil for to expound, then say \emph{atāvatu} again
\end{itemize}

\begin{tcolorbox}[breakable,colback=teal!4,colframe=teal!35!black,title={Before you move on}]
What does \ta{அதாவது} mean? (\textquotedblleft{}That is to say; in other words\textquotedblright{}.) Say it once more.
\end{tcolorbox}
\section[\texorpdfstring{\ta{எடுத்துக்காட்டு} (eṭuttukkāṭṭu)}{எடுத்துக்காட்டு (eṭuttukkāṭṭu)}]{\ta{எடுத்துக்காட்டு} (eṭuttukkāṭṭu) — an example}

\label{lesson:TA-C351-etuttukkattu}



\begin{quote}
Before the new one: say the Tamil for to expound, then the Tamil for that is to say.
\end{quote}

\subsection*{You'll want to know: \ta{எடுத்துக்காட்டு}}
\textbf{\ta{எடுத்துக்காட்டு}} — \emph{eṭuttukkāṭṭu} — \textquotedblleft{}an example\textquotedblright{}.

\begin{quote}
\ta{ஆசிரியர்} \ta{ஒரு} \ta{எளிய} \ta{எடுத்துக்காட்டு} \ta{கொடுத்தார்}.
\end{quote}

\begin{quote}
\emph{āciriyar oru eḷiya eṭuttukkāṭṭu koṭuttār.}
\end{quote}

\begin{quote}
The teacher gave a simple example.
\end{quote}

The pure-Tamil word beside \ta{உதாரணம்}: \ta{எடுத்துக்காட்டாக}, for example.

\begin{cousinweb}
\textbf{\ta{எடுத்துக்காட்டு}} (\emph{eṭuttukkāṭṭu}) is native: \emph{eṭuttu}, \textquotedblleft{}having taken up\textquotedblright{} (DEDR 851), and \emph{kāṭṭu}, \textquotedblleft{}a showing\textquotedblright{} (DEDR 1443): something picked up and shown.
\end{cousinweb}

\begin{grammarlens}[title={What you've built}]
One more word for giving reasons and opinions.
\end{grammarlens}

\subsection*{Your turn}
\begin{itemize}
\raggedright
  \item \emph{Say it:} \emph{eṭuttukkāṭṭu}
  \item \emph{Say it:} \emph{eṭuttukkāṭṭu}, once more
  \item \emph{Say it:} \emph{āciriyar oru eḷiya eṭuttukkāṭṭu koṭuttār.} — The teacher gave a simple example.
  \item \emph{Recall:} say the Tamil for to expound, then the Tamil for that is to say, then say \emph{eṭuttukkāṭṭu} again
\end{itemize}

\begin{tcolorbox}[breakable,colback=teal!4,colframe=teal!35!black,title={Before you move on}]
What does \ta{எடுத்துக்காட்டு} mean? (\textquotedblleft{}An example\textquotedblright{}.) Say it once more.
\end{tcolorbox}
